% ============================================================================
% IV. Процедура и доказателствена стойност
% ============================================================================
\section{Процедура и доказателствена стойност}

Възстановеният видеозапис има стойност за разследването само ако съдът може да се
довери, че той отразява съдържанието на иззетия носител и не е създаден или изменен
при самата експертиза. Затова техническите методи от раздел~III трябва да се прилагат
в рамките на процедура, която може да бъде проверена. Стандартът ISO/IEC 27037 дава
насоки за идентифициране, събиране, придобиване и съхраняване на цифрови доказателства
\cite{iso27037}. Ръководството на NIST SP 800-86 описва процеса на цифровата
криминалистика в четири етапа: събиране, изследване, анализ и докладване
\cite{nist80086}.

\subsection{Работа върху копие}

Първото правило е, че оригиналният носител не се променя. Възстановяването се извършва
върху побитово копие (образ) на носителя, а не върху самия носител
\cite{nist80086}. Това е особено важно при възстановяване, защото всяка
операция за запис върху оригинала, дори случайна, може да презапише точно онези
освободени клъстери, които съдържат изтрития видеозапис.

\subsection{Контролни суми и верига на съхранение}

Целостта на образа се доказва с криптографска контролна сума (хеш), изчислена при
придобиването. Повторното изчисляване по всяко време трябва да даде същия резултат. Това
позволява на независим експерт да провери, че работата е извършена върху същите данни
\cite{nist80086}. Всяко действие с носителя и с образа се документира: кой, кога и с
каква цел го е извършил. Тази документация образува \term{веригата на съхранение}
(chain of custody) \cite{nist80086}.

Контролната сума се изчислява и за всеки възстановен файл. Така в доклада може да се
посочи точно кой файл е обект на последващия видеоанализ и къде в образа на носителя
се намират данните, от които е получен.

\subsection{Възстановено и реконструирано съдържание}

Методите от раздел~III се различават по това колко от резултата идва от носителя и
колко е добавено при експертизата. Тази разлика трябва да е видима в доклада.

\begin{itemize}[leftmargin=2cm]
  \item При възстановяване по метаданни на файловата система и при изрязване на
        непрекъснат файл резултатът е копие на байтове, които физически са на носителя.
        Твърдението е проверимо: всеки байт може да бъде намерен в образа.
  \item При сглобяване на фрагментиран файл байтовете също са от носителя, но
        тяхната подредба е решение на експерта или на програмата. Трябва да се посочи кои
        участъци са съединени и как е проверено, че те принадлежат на един файл.
  \item При ремонт с еталонен файл новата кутия \code{moov} не идва от носителя. Тя е
        построена от програмата въз основа на друг файл. Видеоданните са оригинални, но
        тяхното разделяне на кадри и времето на възпроизвеждане са изчислени. Това трябва
        да се отбележи изрично, заедно с произхода на еталонния файл.
  \item При възстановяване на ниво кадри се възстановяват отделни кадри, а не файлът.
        Липсващите кадри и участъци трябва да се посочат, за да не се създава впечатление
        за непрекъснат запис.
\end{itemize}

Същото важи и за времевите данни. Времето на всеки кадър във файл MP4 се изчислява от
таблиците в \code{moov} \cite{iso14496_12}. Ако \code{moov} е построена наново,
изчисленото време на кадрите е резултат от реконструкцията, а не оригинален запис.
Изводите за продължителност на събитие или за поредност на действия трябва да отчитат
това.

\subsection{Ограничения}

Методите имат и общи граници, които експертът трябва да заяви. Презаписаните данни са
загубени и никой метод не ги връща \cite{na2014}. При твърдотелни дискове освободените
блокове може да са изчистени от самия носител, така че изрязването да не намери нищо
\cite{bell2010}. Съвременните мобилни телефони освен това често криптират данните, и
без ключа възстановените байтове не могат да бъдат прочетени. Липсата на кадри не може да се тълкува
като липса на събития в записания интервал. Накрая, успешното декодиране доказва само,
че данните са валиден видеопоток. То не доказва, че записът не е бил редактиран преди
изтриването. Проверката за автентичност на съдържанието е отделна задача на
видеоанализа и излиза извън обхвата на тази работа.
